\section{为什么要从 RLHF 转到 RLVR}

从工程和研究两个角度看，人类反馈都很难规模化。它贵、慢、噪声大，而且容易被 hack。于是课堂自然把问题转到另一个方向：\textbf{有没有一种奖励，我们能直接验证它对不对？}

这就是 \textbf{verifiable rewards} 的出发点。数学、代码、某些形式化推理题，都是典型例子。只要任务有可验证的终局 reward，我们就可以把 RL 真正跑起来，而不是只盯着人类主观偏好。

\begin{figure}[H]
\centering
\includegraphics[width=0.95\textwidth]{slides-images/slide-014.jpg}
\caption{前半段训练把模型推到类 GPT-3.5，今天的目标是继续把它往 o1 / r1 的风格推进}
\end{figure}

\begin{figure}[H]
\centering
\includegraphics[width=0.95\textwidth]{slides-images/slide-015.jpg}
\caption{核心问题：如果 RLHF 会过优化，那能不能切到 reward 更明确、可精确验证的领域}
\end{figure}

\begin{figure}[H]
\centering
\includegraphics[width=0.95\textwidth]{slides-images/slide-016.jpg}
\caption{本讲结构：先讲算法，再讲案例；算法部分从 PPO 走到 GRPO}
\end{figure}

\begin{knowledgebox}{为什么 RLVR 适合当作 post-training 的试验田}
这类任务最吸引人的地方，不只是 reward 更干净，而是它让研究者可以把主要精力放在\textbf{算法、采样、搜索与 credit assignment} 上，而不用先解决“人类是否喜欢这个回答”的测量难题。也正因为如此，数学、代码、可执行工具调用常被用来观察 RL 是否真的把模型推向了更强的推理行为，而不是仅仅学会迎合评测器。
\end{knowledgebox}

